% 爱德华·康登（综述）
% license CCBYSA3
% type Wiki

本文根据 CC-BY-SA 协议转载翻译自维基百科\href{https://en.wikipedia.org/wiki/Edward_Condon}{相关文章}
\begin{figure}[ht]
\centering
\includegraphics[width=6cm]{./figures/748cfbd4f06fd931.png}
\caption{} \label{fig_ADHkd_1}
\end{figure}
爱德华·尤勒·康登（Edward Uhler Condon，1902年3月2日－1974年3月26日）是一位美国核物理学家、量子力学的先驱之一，并在第二次世界大战期间参与了雷达的研制工作，以及在曼哈顿计划中短暂参与了核武器的开发。以他命名的“弗兰克–康登原理”和“斯莱特–康登规则”均纪念了他在该领域的贡献。\(^\text{[1][2][3]}\)

他曾于1945年至1951年间担任美国国家标准局（现为美国国家标准与技术研究院，NIST）的第四任局长。1946年，康登出任美国物理学会会长；1953年，他担任美国科学促进会会长。

在麦卡锡主义时期，康登是最早成为美国众议院非美活动调查委员会攻击目标的知名科学家之一。1948年，他被公开指控为“美国原子安全体系中最薄弱的一环”，理由是他掌握大量机密信息、与原子弹研制有关联，并被指涉嫌同情共产主义和苏联。康登事件成为当时反对麦卡锡主义者（尤其是科学界）关注的标志性案件之一，舆论广泛关注，他也得到了包括美国总统哈里·杜鲁门在内的众多知名科学家的公开支持。\(^\text{[4]}\)

1968年，康登作为《康登报告》的主要作者再次广为人知。该报告是由美国空军资助的官方评估，得出的结论是“不明飞行物都有平凡的解释”。月球上的“康登环形山”即以他的名字命名，以纪念他的贡献。
\subsection{背景}
\begin{figure}[ht]
\centering
\includegraphics[width=6cm]{./figures/3223424e58a1eb6d.png}
\caption{图1. 弗兰克–康登原理能量图。由于电子跃迁相比于原子核运动要快得多，因此在跃迁过程中，最有利的振动能级是那些对应于原子核坐标变化最小的能级。图中所示的势阱表明，跃迁在 $v = 0$与$v = 2$之间最为有利。} \label{fig_ADHkd_2}
\end{figure}
爱德华·尤勒·康登于1902年3月2日出生在新墨西哥州的阿拉莫戈多，父母分别是威廉·爱德华·康登和卡罗琳·尤勒。他的父亲当时正在监督一条窄轨铁路的建设，\(^\text{[5][6]}\)当地的许多铁路都是由伐木公司修建的。1918年，康登在加利福尼亚州奥克兰完成高中学业后，曾在《奥克兰询问报》及其他报纸担任记者三年。\(^\text{[5]}\)他具有爱尔兰血统。\(^\text{[7]}\)

随后，他进入加利福尼亚大学伯克利分校就读，最初加入化学院；后来得知他高中时的物理老师成为了该校教师后，他转而改修理论物理。\(^\text{[8]}\)康登仅用三年时间获得学士学位，又用两年取得博士学位。\(^\text{[5]}\)他的博士论文结合了雷蒙德·塞耶·伯奇在带状光谱强度测量与分析方面的研究，以及詹姆斯·弗兰克的一个建议。\(^\text{[9][10]}\)

凭借国家研究委员会奖学金，康登前往德国哥廷根大学师从马克斯·玻恩，又赴慕尼黑大学师从阿诺德·索末菲。在后者指导下，康登以量子力学的形式重写了自己的博士论文，并提出了著名的“弗兰克–康登原理”。\(^\text{[9]}\)

1927年秋，他在《物理评论》上看到一则招聘广告后，加入贝尔电话实验室从事公共关系工作，主要负责宣传他们关于电子衍射的发现。\(^\text{[5][11]}\)在贝尔实验室工作期间，他还研究了词语使用频率的分布，并将结果与心理学中的韦伯–费希纳定律联系起来。\(^\text{[12]}\)
\subsection{职业生涯}
\subsubsection{早期经历}
\begin{figure}[ht]
\centering
\includegraphics[width=6cm]{./figures/6f359ea6e05a3107.png}
\caption{阿诺德·索末菲，德国理论物理学家，曾在20世纪20年代指导过康登。} \label{fig_ADHkd_3}
\end{figure}
康登曾在哥伦比亚大学短暂任教，并于1928年至1937年期间担任普林斯顿大学物理学副教授，\(^\text{[5]}\)其间曾在明尼苏达大学任教一年。\(^\text{[13]}\)他与菲利普·M·莫尔斯合著了《量子力学》，这是1929年出版的第一本英文版量子力学教材。随后，他与G.H.肖特利合著了《原子光谱理论》，该书自1935年出版起便被誉为该领域的“权威圣经”。\(^\text{[14][15][16][Note 1]}\)

1937年起，康登担任位于匹兹堡的西屋电气公司研究副主任，在那里他创建了核物理、固体物理和质谱学的研究项目。随后，他领导公司在微波雷达开发方面的研究工作。\(^\text{[13]}\)他还参与了用于分离原子弹原料铀的设备研发。\(^\text{[5]}\)在第二次世界大战期间，康登担任国家防务研究委员会的顾问，并协助组织了麻省理工学院的辐射实验室。\(^\text{[13][15]}\)1940年5月11日，康登在纽约世界博览会上展示了名为“Nimatron”的机电装置，这是一台用于玩尼姆游戏的早期计算机式设备。\(^\text{[18]}\)
\subsubsection{政府任职}
\begin{figure}[ht]
\centering
\includegraphics[width=6cm]{./figures/1ec8df323ef9ca2e.png}
\caption{} \label{fig_ADHkd_4}
\end{figure}